\section{Strict retractions and prescribed fence homotopies}
\label{sec:D}

All categories in this section are $\mathcal U$-small. A categorical homotopy is a finite zigzag of natural transformations. We write $Z$ for the free category on a finite oriented interval, and distinguish its initial vertex $0$ and final vertex $z$ (these names specify positions, not initial or terminal objects in the category). A padding of $Z$ is a longer fence $Z'$ with a functor $Z'\to Z$ preserving these endpoint positions; added edges may map to identities. We use the model structure of \citet[Theorem~4.9]{D_Thomason1980} and the corrected fact that every Thomason cofibration is a full sieve and a pseudo-Dwyer map \citep[Proposition~2]{D_Cisinski1999}.

\begin{claim}{D.1}{Proved}{Strict retraction in the prescribed weak-homotopy class}
Let $A$ be Thomason-fibrant, let $i:A\hookrightarrow B$ be a Thomason cofibration, and let $p:B\to A$ satisfy $pi\simeq\mathrm{id}_A$ by a finite categorical homotopy. There is $r:B\to A$ with $ri=\mathrm{id}_A$ in $\mathrm{Cat}$ and $[r]=[p]$ in $\mathrm{Hot}$.
\end{claim}
\begin{proof}
Represent the given homotopy by $H:A\times Z\to A$, with $H(-,0)=pi$ and $H(-,z)=\mathrm{id}_A$. The constant-path functor $c:A\to\mathrm{Fun}(Z,A)$ is a categorical homotopy equivalence: delete the endpoints of the fence successively, using each incident arrow for the natural transformation between the identity and the resulting endpoint-collapse endofunctor, and compose the resulting finite zigzags. Consequently $c$ is a Thomason weak equivalence.

Factor the endpoint evaluation functor as
\[
 \mathrm{Fun}(Z,A)\xrightarrow{j}P
 \xrightarrow{(d_0,d_1)}A\times A,
\]
with $j$ a trivial cofibration and the second map a fibration. The maps $d_0,d_1:P\to A$ are fibrations because the projections $A\times A\to A$ are pullbacks of the fibration $A\to *$. The composite $jc$ is a weak equivalence and $d_0jc=d_1jc=\mathrm{id}_A$, so both $d_0,d_1$ are weak equivalences as well. Thus both are trivial fibrations. Let $h:A\to P$ be $j$ applied to the adjoint of $H$. The square with top $h$, left $i$, bottom $p$, and right $d_0$ commutes. Lift it to $\widetilde h:B\to P$ and set $r=d_1\widetilde h$. The equations $ri=d_1h=\mathrm{id}_A$ and $d_0\widetilde h=p$ hold in $\mathrm{Cat}$. In $\mathrm{Hot}$ both $d_0$ and $d_1$ equal $(jc)^{-1}$; hence $[r]=[p]$ there.
\end{proof}

This confirms every step of the weaker statement in CATHOT E.6. The factorization supplies no finite natural-transformation zigzag on $B$ between $r$ and $p$.

\begin{claim}{D.2}{Proved}{The local-target and groupoid cases}
Under the hypotheses of D.1, a local target $A$ admits $r$ with $[r]=[p]$ in $\mathrm{Cat}_{\mathrm{hot}}$. For an actual groupoid target the same conclusion follows for every full inclusion $i$ that is injective on objects, without Thomason fibrancy or cofibration assumptions.
\end{claim}
\begin{proof}
For the first assertion, locality makes the comparison
$\pi_0\mathrm{Fun}(B,A)\to[BB,BA]$ injective, and D.1 supplies equal images.
For the second, the given zigzag composes, after reversing its invertible components where necessary, to a natural isomorphism $\alpha:pi\Rightarrow\mathrm{id}_A$. Identify $A$ with its full image in $B$. For $b\in A$ put $r(b)=b$ and $\beta_b=\alpha_b:p(b)\to b$; for $b\notin A$ put $r(b)=p(b)$ and $\beta_b=\mathrm{id}$. For $u:b\to b'$ define
$r(u)=\beta_{b'}p(u)\beta_b^{-1}$.
The identities and composition equations follow by cancellation of consecutive $\beta_b^{-1}\beta_b$. Naturality of $\alpha$ proves $ri=\mathrm{id}_A$ on all arrows as well as objects. The family $\beta$ is a natural isomorphism $p\Rightarrow r$.
\end{proof}

\begin{claim}{D.3}{Proved}{Prescribed fence extension for targets with a terminal object}
Let $i:C\hookrightarrow D$ be a full sieve inclusion, let $F$ have a terminal object $t$, and let $f:D\to F$ and $H:C\times Z\to F$ agree on $C\times\{0\}$. There is an extension over $D\times Z'$ for a padding $Z'\to Z$ obtained by adding one edge at the initial position.
\end{claim}
\begin{proof}
Define $f':D\to F$ by $f'|_C=f|_C$ and $f'(d)=t$ outside $C$. Since $C$ is a sieve, no arrow has source outside $C$ and target in $C$. Every remaining mixed arrow maps to the unique arrow into $t$. The identity components on $C$ and the terminal arrows outside $C$ define $f\Rightarrow f'$. Next define $\widehat H:D\times Z\to F$ by $H$ on $C\times Z$ and the constant value $t$ outside. Mixed naturality squares commute by uniqueness of arrows into $t$. Prepend the transformation $f\Rightarrow f'$ and map the added edge to the identity at $0\in Z$. Its restriction to $C$ is stationary, so the resulting extension has exactly the required boundary.
\end{proof}

\begin{claim}{D.4}{Proved}{Prescribed fence extension for targets with an initial object}
Let $i:C\hookrightarrow D$ be a pseudo-Dwyer inclusion. Choose its cosieve neighborhood $W\subseteq D$, a retraction $s:W\to C$, and a transformation $\varepsilon:is\Rightarrow\mathrm{id}_W$ whose restriction to $C$ is the identity. Let $F$ have an initial object $o$. Every pair $f:D\to F$, $H:C\times Z\to F$ agreeing at the initial endpoint admits the extension of D.3 after one padding edge. When $W=D$, the conclusion holds for every target $F$.
\end{claim}
\begin{proof}
On $W$ put $f'=fis$, and outside $W$ put $f'=o$. There are no arrows from $W$ to its complement because $W$ is a cosieve. Define arrows with source outside $W$ and target in $W$ using the unique maps from $o$. This specifies a functor, and $f\varepsilon$ on $W$, together with the initial arrows elsewhere, specifies $f'\Rightarrow f$. Its restriction to $C$ is the identity. Extend $H$ to $D\times Z$ by $H(s(-),-)$ on $W$ and the constant initial object on the complement. The same cosieve and uniqueness checks prove functoriality and naturality. Prepend the reversed-in-position edge $f\Leftarrow f'$, mapped to the identity at $0\in Z$. This is the desired padding. With $W=D$ the construction has no complementary part and uses no initial object.
\end{proof}

\begin{claim}{D.5}{Proved}{A non-cofibrant corner can still admit padded extension}
Let $C=\{0\}$, $D=[1]$, and let $F$ be the three-object poset with arrows $o\to u$ and $o\to v$. The corner functor sending $(0,0)$ to $o$, $(1,0)$ to $u$, and $(0,1)$ to $v$ has no extension to $D\times[1]$. It extends over a padded fence with two edges. The category $F$ is Thomason-fibrant.
\end{claim}
\begin{proof}
An extension over the square requires an upper bound of $u,v$, which does not exist in $F$. Over the fence $0\leftarrow0'\to1$, mapping its first edge to $\mathrm{id}_0$ in $[1]$, take the successive functors $D\to F$ to be $(o\to u)$, $(o\xrightarrow{\mathrm{id}}o)$, and $(v\xrightarrow{\mathrm{id}}v)$. The two transformations are $f_0\Leftarrow f_{0'}\Rightarrow f_1$, and restriction to $C$ is $o\Leftarrow o\Rightarrow v$, the required pullback of the original boundary.
To check fibrancy, apply the construction underlying D.4 to each generating Thomason horn inclusion: it is Dwyer, so its cosieve neighborhood has a right-adjoint retraction. Given a map from the horn to $F$, define its extension by that retraction on the neighborhood and by $o$ elsewhere. The same arrow check proves a strict extension. Thus $F\to *$ has the required right lifting property.
\end{proof}


\begin{claim}{D.6}{Refuted}{Fibrancy does not make the categorical comparison injective}
There are a finite category $X$ with six objects and twenty-two arrows, a $\mathcal U$-small Thomason-fibrant poset $F$, and two functors $u,v:X\to F$ such that $[u]\ne[v]$ in $\mathrm{Cat}_{\mathrm{hot}}$, while $Bu$ and $Bv$ are unpointedly homotopic. In particular fibrancy of the target does not imply injectivity of
\[
 \pi_0\mathrm{Fun}(X,F)\longrightarrow[BX,BF].
\]
\end{claim}
\begin{proof}
Let $P$ have two elements in each of three levels $0,1,2$, with $x<y$ exactly when the level of $x$ is smaller than the level of $y$. Its order complex is the simplicial join of three copies of $S^0$, so $BP$ is homeomorphic to $S^2$. Let $X$ be the free category on the Hasse quiver of $P$: there are four arrows from level $0$ to level $1$, four from level $1$ to level $2$, and eight length-two composites. Together with the six identities these are all twenty-two arrows. Let $q:X\to P$ send a path to its endpoints.

For every poset $T$, precomposition gives an isomorphism of categories
\[
 q^*:\mathrm{Fun}(P,T)\xrightarrow{\cong}\mathrm{Fun}(X,T).
\]
Indeed a functor on the free quiver specifies an order-preserving assignment on the eight generating comparisons. Transitivity forces exactly the four additional level-$0$ to level-$2$ comparisons; uniqueness of arrows in $T$ identifies the two composites between each such pair. A natural transformation on either side is exactly a pointwise comparison of the same object assignments, and the naturality squares commute because $T$ is a poset. This proves the claimed isomorphism on both objects and arrows.

Take a fibrant replacement $j:P\to F$ in the Thomason model structure on posets, whose existence is \citet[Theorem~2.6]{D_Raptis2010}. The inclusion of posets into categories creates fibrations and weak equivalences, so $F$ is Thomason-fibrant as a category and $Bj$ is a weak equivalence. Set $u=jq$ and let $v$ be any constant functor. An equality $[u]=[v]$ in $\mathrm{Cat}_{\mathrm{hot}}$ would, by the displayed isomorphism, give a finite natural-transformation zigzag from $j:P\to F$ to a constant functor. Applying the nerve and realization gives a homotopy between $Bj$ and a constant, contradicting the isomorphism $H_2(Bj;\mathbb Z):\mathbb Z\to\mathbb Z$.

On the other hand, $NX$ collapses onto the Hasse graph. Each of its eight nondegenerate two-simplices has its own composite long edge; that edge belongs to no other nondegenerate two-simplex. Delete each triangle together with its free long edge. The resulting connected graph has six vertices and eight edges, hence is a wedge of three circles. The space $BF$ is simply connected because it is weakly equivalent to $S^2$. Therefore every map from this graph to $BF$ is nullhomotopic: contract a spanning tree, and nullhomotope each of its three loop maps using $\pi_1BF=0$. It follows that $Bu\simeq Bv$ as claimed. All category compositions, free-face incidences, and the chain-complex ranks are checked in \texttt{code/d\_finite\_checks.py}; the infinite fibrant replacement is justified by the specified model structure, not by a bounded computation.
\end{proof}


\begin{claim}{D.7}{Proved}{Backward transformations extend along sieves}
Let $i:C\hookrightarrow D$ be a full sieve, let $f:D\to F$, and let $\alpha:g\Rightarrow fi$ be a natural transformation. There is $\widetilde g:D\to F$ with $\widetilde g i=g$ and a transformation $\widetilde\alpha:\widetilde g\Rightarrow f$ restricting to $\alpha$. Its components outside $C$ are identities.
\end{claim}
\begin{proof}
Set $\widetilde g=g$ on $C$ and $\widetilde g=f$ on the full complement. For a mixed arrow $u:c\to d$, set $\widetilde g(u)=f(u)\alpha_c$. There are no arrows in the other mixed direction. Naturality of $\alpha$ checks compositions whose first two objects lie in $C$; functoriality of $f$ checks compositions whose last two objects lie outside. These are all mixed cases. The components $\alpha_c$ on $C$ and identities elsewhere give the required transformation.
\end{proof}

\begin{claim}{D.8}{Proved}{An exact reduction to forward endpoint extension}
The universal strict-retraction assertion of Part D is equivalent to the following endpoint-extension assertion: for every Thomason cofibration $j:C\hookrightarrow D$, every Thomason-fibrant $F$, every $f:D\to F$, and every transformation $\alpha:fj\Rightarrow h:C\to F$, there is $g:D\to F$ with $gj=h$ and $g\simeq f$ by a finite categorical homotopy. This equivalence makes no assertion that either universal property holds.
\end{claim}
\begin{proof}
First construct a particular category from the endpoint-extension data. Put $h_0=fj$ and $h_1=h$. Let $A$ be the collage of $h_1$: its objects are the disjoint union of the objects of $C$ and $F$, its full subcategories on these sets are $C$ and $F$, and
\[
 A(c,y)=F(h_1(c),y),\qquad A(y,c)=\varnothing
 \quad(c\in C,\ y\in F).
\]
Composition uses composition in $C,F$ and the functor $h_1$. There are a reflector $q:A\to F$, equal to $h_1$ on $C$ and the identity on $F$, its fully faithful section $k:F\to A$, and a unit $\eta:\mathrm{id}_A\Rightarrow kq$ with identity components on $F$.

The category $A$ is Thomason-fibrant. To verify this directly, consider a generating horn inclusion $K\hookrightarrow L$ and $v:K\to A$. This inclusion is a full sieve. Fibrancy of $F$ extends $qv$ to $w:L\to F$. On $K$ retain $v$, and on objects outside $K$ use $kw$. For a mixed arrow $x\to y$ use the composite
$ v(x)\xrightarrow{\eta_{v(x)}}kqv(x)\xrightarrow{kw(x\to y)}kw(y)$.
The unit naturality equation checks all compositions crossing the sieve, exactly as in D.7. Thus every generating horn lifts.

Define $e:A\to A$ to be $kh_0$ on $C$, the identity on $F$, and to send a mixed arrow $\beta:h_1(c)\to y$ to $\beta\alpha_c:h_0(c)\to y$. Naturality of $\alpha$ proves functoriality. There are transformations $e\Rightarrow kq$ (components $\alpha$ on $C$) and $\mathrm{id}_A\Rightarrow kq$ (the unit). Form
$B=A\amalg_C D$. Then $i:A\hookrightarrow B$ is a Thomason cofibration. The functors $e$ on $A$ and $kf$ on $D$ agree on $C$, giving $p:B\to A$ with $pi=e\simeq\mathrm{id}_A$.

A strict retraction $r:B\to A$ in the finite categorical homotopy class of $p$ gives $g=qr|_D$. Its restriction to $C$ is $q|_C=h_1$, and $g\simeq qk f=f$. This proves that the universal strict-retraction property implies the stated endpoint property. Conversely apply the endpoint property successively to every forward edge of a finite zigzag on a cofibration's source, and use D.7 for every backward edge. At each stage the extended endpoint remains categorically homotopic to the preceding one. Applied to $pi\simeq\mathrm{id}_A$, this gives the required strict retraction. These constructions prove both implications.
\end{proof}

Write $P_4=c\operatorname{Sd}^2\Delta^1$, with its five vertices in the order
\[
 0\longrightarrow1\longleftarrow2\longrightarrow3\longleftarrow4.
\]
For objects $a,b$ of a category $F$, let $\mathscr P_F(a,b)$ be the category of functors $P_4\to F$ taking the endpoints to $a,b$, with natural transformations that are identities at both endpoints.

\begin{claim}{D.9}{Proved}{A uniform eight-step comparison of four-edge paths}
For every Thomason-fibrant $F$, realization induces a bijection
\[
 \pi_0\mathscr P_F(a,b)\ \xrightarrow{\ \cong\ }\ \Pi_1(BF)(a,b).
\]
Two objects representing the same fundamental path class are connected by eight natural transformations, with orientations
$\Rightarrow,\Leftarrow,\Rightarrow,\Leftarrow,\Rightarrow,\Leftarrow,\Rightarrow,\Leftarrow$,
all identity at the endpoints. The bound does not depend on $F,a,b$, or the paths.
\end{claim}
\begin{proof}
Put $S=\operatorname{Ex}^2 NF$, which is Kan by Thomason fibrancy. Its vertices are objects of $F$, and its edges are exactly the displayed functors $P_4\to F$. The adjunction counit $\operatorname{Sd}^2S\to NF$, together with the subdivision homeomorphism, identifies fundamental path classes of $S$ with those of $BF$ and sends an edge to its four-edge realization.

We recall the elementary Kan edge calculation, including the feature needed here. Every fundamental path class in a Kan complex is represented by a single edge: the two-dimensional horn fillers compose adjacent edges and reverse an edge. Two edges $a,b:x\to y$ represent the same class precisely when there is a two-simplex with edges $01=a$, $02=b$, and $12=\mathrm{id}_y$. For completeness, the relation defined by such a simplex is reflexive by $s_1a$. Symmetry follows by filling the three-dimensional horn with faces $012$ the given simplex, $013=s_1a$, and $123$ constant; its fourth face reverses the relation. Transitivity follows by filling the horn with faces $012$ and $023$ the two given relation simplices and $123$ constant. Two-dimensional horns give composition of the resulting edge classes. Three-dimensional horns give independence of the chosen composites and associativity; a two-dimensional horn with one constant edge gives inverses. Hence these classes form a groupoid with precisely the relations $[02]=[12][01]$ of all two-simplices. The realization's two-skeleton has one edge for every nondegenerate one-simplex and one triangular relation for every nondegenerate two-simplex; cells of dimension at least three do not change its fundamental groupoid. This proves the asserted calculation.

Let $p,q:P_4\to F$ have the same fundamental path class. Take the resulting two-simplex of $S$. By adjunction it is a functor
\[
 \sigma:Q=c\operatorname{Sd}^2\Delta^2\longrightarrow F
\]
whose edges $01,02$ are $p,q$ and whose whole edge $12$ is the constant functor at $b$. We now give the universal eight-step comparison inside $Q$.

The objects of $Q$ are nonempty chains of nonempty subsets of $\{0,1,2\}$, ordered by inclusion of chains. Write $T=012$; thus $[0,01,T]$ denotes the chain $\{0\}\subset\{0,1\}\subset\{0,1,2\}$. The following rows specify nine functors $P_4\to Q$:
\[
\begin{array}{c|ccccc}
 &0&1&2&3&4\\\hline
 R_0 &[0]&[0,01]&[01]&[1,01]&[1]\\
 R_1 &[0]&[0,01,T]&[01,T]&[1,01,T]&[1]\\
 R_2 &[0]&[0,T]&[T]&[1,T]&[1]\\
 R_3 &[0]&[0,T]&[T]&[1,12,T]&[1,12]\\
 R_4 &[0]&[0,T]&[T]&[12,T]&[12]\\
 R_5 &[0]&[0,T]&[T]&[2,12,T]&[2,12]\\
 R_6 &[0]&[0,T]&[T]&[2,T]&[2]\\
 R_7 &[0]&[0,02,T]&[02,T]&[2,02,T]&[2]\\
 R_8 &[0]&[0,02]&[02]&[2,02]&[2]
\end{array}
\]
Each row has its first, third, and fifth entries included in their adjacent entries, as required by the arrows of $P_4$. Inclusion of chains also gives
$R_0\le R_1\ge R_2\le R_3\ge R_4\le R_5\ge R_6\le R_7\ge R_8$ pointwise. Its first endpoint is always $[0]$. Its second endpoint travels entirely along the subdivided edge $12$. Composing with $\sigma$ therefore yields the claimed eight transformations, identity at both endpoints, from $p$ to $q$. Conversely any endpoint-fixed natural transformation realizes to an endpoint-fixed path homotopy. This establishes injectivity; the Kan edge calculation gives surjectivity.

The program \texttt{code/d\_prism\_probe.py} checks every entry and comparison in this table inside the actual twenty-five-object poset $Q$. As an additional finite check it collapses the five-object edge $12$, enumerates all thirty-one endpoint-fixed $P_4$ paths in the resulting twenty-one-object poset, and finds comparison distance eight between the two specified paths. The proof uses the lifted table, so it requires no identification of a categorical quotient with a poset quotient.
\end{proof}

\begin{claim}{D.10}{Proved}{Full prescribed extension for a family of one-cell attachments}
Let $F$ be Thomason-fibrant, and form the pushout
\[
 D=C\amalg_{\coprod_{s\in J}\{0,4\}}\coprod_{s\in J}P_4
\]
for any $\mathcal U$-small set $J$ and any endpoint maps into $C$. For every $f:D\to F$ and prescribed finite fence homotopy $H:C\times Z\to F$ beginning at $f|_C$, there is an extension $D\times Z'\to F$ over a padding $Z'\to Z$. A fence with $L$ edges needs at most $9L$ edges after padding. Consequently the strict-retraction assertion of Part D holds for every such attachment to a Thomason-fibrant $C$, including the single attachment to fibrant poset models of $S^1$ or $S^2$ requested in the prompt.
\end{claim}
\begin{proof}
It suffices to extend one prescribed transformation. A backward transformation extends along every sieve by D.7. For a forward transformation $\alpha:h_0=f|_C\Rightarrow h_1$, examine one cell. Write $p:P_4\to F$ for its original map, $x_0,x_1$ for its endpoints, and $y_0,y_1$ for the new endpoints. The endpoint components are arrows $\alpha_j:x_j\to y_j$.

By D.9 choose a four-edge path $p':y_0\leadsto y_1$ in the fundamental class
$[\alpha_1][p][\alpha_0]^{-1}$, where composition is read from right to left. Define $q:P_4\to F$ to retain the three interior objects and the two interior arrows of $p'$, to have endpoints $x_0,x_1$, and to send the two exterior arrows to
\[
 p'(0\to1)\alpha_0,
 \qquad p'(4\to3)\alpha_1.
\]
There is a natural transformation $q\Rightarrow p'$ with endpoint components $\alpha_0,\alpha_1$ and identity interior components. Moreover
$[q]=[\alpha_1]^{-1}[p'][\alpha_0]=[p]$.
D.9 supplies eight transformations from $p$ to $q$, identity at the endpoints and with its fixed alternating orientation pattern. Use the constant functor $h_0$ on $C$ through these eight steps, and then use $\alpha$ on $C$ and $q\Rightarrow p'$ on the cell. Their endpoint restrictions agree, so the pushout property gives the required functors and transformations on $D$.

For an arbitrary set of cells, make the same choices for each cell. The uniform eight-step orientation pattern lets all cells be treated simultaneously, even for an infinite $J$. Products with a fixed category preserve colimits in $\mathrm{Cat}$ because they are left adjoints to functor categories, so this gluing is equivalently a functor on $D$ times the resulting fence. The first eight edges map to the identity at the old boundary vertex, and the ninth maps to the given forward edge. Repeat the procedure along the prescribed fence, using one edge for a backward transformation and nine for a forward transformation. For the retraction assertion take $F=C$ and apply this extension to the given homotopy $pi\simeq\mathrm{id}_C$.
\end{proof}


\begin{claim}{D.11}{Proved}{Relative comparison in every fixed simplex dimension}
Put $P_n=c\operatorname{Sd}^2\Delta^n$ and $\partial P_n=c\operatorname{Sd}^2\partial\Delta^n$. Let $F$ be Thomason-fibrant and let $p,q:P_n\to F$ agree on $\partial P_n$. Their adjoint simplices are vertices of the fiber of the Kan fibration
\[
 (\operatorname{Ex}^2NF)^{\Delta^n}\longrightarrow
 (\operatorname{Ex}^2NF)^{\partial\Delta^n}
\]
over their common boundary. For any two vertices in the same component of this fiber, the functors $p,q$ are connected by $4(n+1)$ natural transformations, all identity on $\partial P_n$, alternating forward and backward starting forward.
\end{claim}
\begin{proof}
Write $S=\operatorname{Ex}^2NF$. We first give a relative cylinder for $\Delta^n$. Let $\pi:[n+1]\to[n]$ identify $n,n+1$ and fix the other vertices. Let $W=\pi^{-1}(\partial\Delta^n)\subset\Delta^{n+1}$. Explicitly, $W$ is the union of the facets missing $i$ for $0\le i<n$ together with their complementary common end facet $V=\Delta^{\{0,\ldots,n-1\}}$ (empty when $n=0$). The latter facet must be included in the gluing. Set
\[
 C_n=\Delta^{n+1}\amalg_W\partial\Delta^n.
\]
The two facets missing $n+1$ and $n$ induce a monomorphism
$\Delta^n\amalg_{\partial\Delta^n}\Delta^n\hookrightarrow C_n$.
To check injectivity, their interiors are disjoint and neither meets $W$; their common facet has vertices $0,\ldots,n-1$; every other face of their boundaries is identified by exactly the prescribed projection to $\partial\Delta^n$. The projection $C_n\to\Delta^n$ has the first facet as a section. The order map on $[n+1]$ sending $n+1$ to $n$ is pointwise at most the identity, preserves $W$, and has the same composite with $\pi$ as the identity. Its simplicial homotopy to the identity descends to $C_n$ and fixes $\partial\Delta^n$. Thus the projection is a simplicial homotopy equivalence. For $n=0$ this construction is simply $C_0=\Delta^1$.

Consequently $C_n$ is a cylinder object of $\Delta^n$ in simplicial sets under $\partial\Delta^n$. That source is cofibrant and $S$ is fibrant. The usual cylinder lifting argument says that maps in the same relative homotopy class are joined using this cylinder: factor its projection as a trivial cofibration followed by a trivial fibration; the first map permits extension into the fibrant target, and lifting the cofibrant source against the latter compares it to any other cylinder. Equivalently, left homotopy for a cofibrant source and fibrant target is independent of the chosen cylinder object \citep[Proposition~A.2.2.1]{CpathsLurieHTT}. Applying this to the given component of the mapping-space fiber gives a map $C_n\to S$ with ends $\widehat p,\widehat q$. Its restriction to $\Delta^{n+1}$ is a simplex whose side faces factor through their prescribed boundary values by $\pi$. Adjunction gives
$\sigma:c\operatorname{Sd}^2\Delta^{n+1}\to F$.

Here is a universal flag comparison between the two end inclusions. A source face is a nonempty subset $T\subseteq\{0,\ldots,n\}$. Start with the assignment $T\mapsto T$. For $k=n+1,n,\ldots,1$, replace each face $T$ with $n\in T$ and $|T|=k$ by $T\cup\{n+1\}$. For each such replacement use two stages: first retain both the old and new image of that face in a flag, and then delete the old image. After these $2(n+1)$ stages, process $k=1,2,\ldots,n+1$ and remove $n$ from the image of each original face of cardinality $k$ containing $n$, again using an intermediate stage retaining both images. This is another $2(n+1)$ stages.

These assignments give functors on the poset of flags. Indeed the transformation of a face depends only on that face, so inclusion of flags is preserved. Each image remains a chain: during insertion, all larger faces have already acquired $n+1$; during deletion, all smaller faces have already lost $n$. A flag contains at most one face of each cardinality, so the intermediate union also remains a chain. Every even stage is contained pointwise in each adjacent odd stage. The resulting $4(n+1)$ comparisons therefore have the stated alternating orientations. The final assignment replaces $n$ by $n+1$, exactly the second end inclusion.

Projecting any intermediate flag by $\pi$ gives precisely the original flag, with duplicate faces deleted. For a boundary flag whose maximal face misses some $i<n$, every intermediate image lies in the corresponding side facet, where $\sigma$ factors through its prescribed boundary by $\pi$. The remaining boundary flags lie in the common facet $\{0,\ldots,n-1\}$ and stay unchanged. Thus applying $\sigma$ makes every boundary component an identity throughout. The formulas prove the assertion in all dimensions. Their chain, functoriality, projection, and boundary conditions are additionally verified through dimension four by \texttt{code/d\_higher\_prism\_checks.py}, checking source posets with up to $1081$ objects.
\end{proof}

\begin{claim}{D.12}{Proved}{Prescribed extension for finite-stage bounded-dimensional cell complexes}
Fix $N\ge1$. A pushout of any $\mathcal U$-small coproduct of the generating Thomason cofibrations
$\partial P_n\hookrightarrow P_n$ with $n\le N$ has the full prescribed-fence extension property into every Thomason-fibrant target. An $L$-edge boundary fence extends over a padding with at most $(4N+5)L$ edges. For a composite of $m$ such stages, with dimension bounds $N_1,\ldots,N_m$, the bound is
\[
 L\prod_{j=1}^{m}(4N_j+5).
\]
Every retract in the arrow category of such a composite has the same extension property. In particular every finite relative Thomason cell complex, and every retract of one, satisfies both full prescribed extension and the strict-retraction assertion of Part D.
\end{claim}
\begin{proof}
Backward transformations are treated by D.7. Consider a prescribed forward transformation $\alpha:h_0\Rightarrow h_1$ on the old category. For one $n$-cell, let $p:P_n\to F$ be the original map, and let $a$ be the edge in
$Y=S^{\partial\Delta^n}$ induced by the boundary transformation, with $S=\operatorname{Ex}^2 NF$. To be explicit about this edge, a transformation is a functor $K\times[1]\to F$. Apply $\operatorname{Ex}^2N$, use preservation of products, precompose the first factor with the unit $\partial\Delta^n\to\operatorname{Ex}^2N(\partial P_n)$, and precompose the second factor with the edge of $\operatorname{Ex}^2N[1]$ corresponding to the functor $P_4\to[1]$ with values $0,0,0,1,1$. This construction is natural under restriction from $P_n$ to $\partial P_n$ and has the prescribed endpoints.

The restriction $X=S^{\Delta^n}\to Y$ is a Kan fibration. Lift $a$ beginning at $\widehat p$ to an edge $u:\widehat p\to\widehat {p'}$. The new cell map $p'$ has boundary $h_1$. Use D.7 on the sieve $\partial P_n\hookrightarrow P_n$ to pull $p'$ back along $\alpha$, obtaining $q:P_n\to F$ with boundary $h_0$ and a transformation $q\Rightarrow p'$ restricting to $\alpha$. The explicit edge construction above gives an edge $v:\widehat q\to\widehat {p'}$ of $X$ over exactly the same edge $a$ of $Y$.

Fill the two-horn in $X$ with vertices $\widehat p,\widehat q,\widehat {p'}$ and edges $02=u$, $12=v$, lying over the degenerate triangle of $Y$ with edges $01=\mathrm{id}_{h_0}$ and $02=12=a$. The fibration supplies the missing edge from $\widehat p$ to $\widehat q$ in the fiber over $h_0$. D.11 now yields $4(n+1)$ natural transformations from $p$ to $q$, identity on the entire cell boundary. Follow them by $q\Rightarrow p'$. On the old category these steps are stationary until the final prescribed transformation $\alpha$.

The orientation pattern is the same in every dimension, and its length is a multiple of four. Add identity comparisons in pairs at the end to reach length $4(N+1)$ for every cell. Zero-cells can stay constant throughout. The cells therefore glue simultaneously, even for an infinite coproduct, giving at most $4N+5$ edges for each original forward edge. Repeat along the original fence. For successive attachment stages, use the previously constructed extension as the prescribed boundary fence of the next stage and compose the padding maps. This multiplies the length bounds.

Finally, write a retract of cofibrations as maps of arrows $i\to j\to i$ whose composite is the identity. Pull the initial map and prescribed boundary homotopy back to $j$ along the retraction of arrows, apply the extension for $j$, and restrict along the section of arrows. The two retraction identities give precisely the original initial map and boundary homotopy, with the same padding. For the strict-retraction assertion, use the extension with target the fibrant source and with prescribed boundary homotopy $pi\simeq\mathrm{id}$.
\end{proof}


\begin{claim}{D.13}{Proved}{No finite cofibration supplies a counterexample}
Let $i:C\hookrightarrow D$ be a Thomason cofibration such that $D$ is finitely presented as a category under $C$. Then $i$ has the full prescribed-fence extension property into every Thomason-fibrant target. In particular this holds for every Thomason cofibration between finite categories. The strict-retraction assertion holds for all such $i$ with fibrant source.
\end{claim}
\begin{proof}
Factor $i$ by the small-object argument as a relative cell map $j:C\to E$ followed by a trivial fibration $q:E\to D$. Since $i$ is a cofibration, there is a lift $s:D\to E$ with $si=j$ and $qs=\mathrm{id}_D$. Every finite set of cells of $E$ is contained in a finite subcomplex under $C$. To see this, each chosen cell has a finite attaching category; each object and arrow of its attaching map is supported by finitely many earlier cells. Choose those dependencies and repeat. The resulting dependency trees are finitely branching and their cell stages strictly decrease along every branch. There is no infinite strictly decreasing sequence of ordinals, so K\"onig's lemma makes each such tree finite. Closing a finite set of cells under these finite dependencies therefore gives a finite subcomplex.

Thus $E$ is the filtered union of its finite relative subcomplexes under $C$. Finite presentation of $D$ under $C$ makes $s$ factor through one of them, say $E_0$. Restricting $q$ gives a retract diagram of arrows between $i$ and the finite relative cell map $C\to E_0$. D.12 applies. A finite category has a finite presentation using its finitely many arrows as generators and its finite composition table as relations. For finite $C,D$, the structural functor from $C$ imposes only finitely many further equations, so $D$ is finitely presented under $C$. This proves the particular assertion as well as the retraction statement.
\end{proof}


\begin{claim}{D.14}{Proved}{Uniform extension for three broad classes of targets}
Every Thomason cofibration has the full prescribed-fence extension property into each of the following targets: a filtered category; a category with all pushouts; a category with all pullbacks. For filtered targets and targets with all pushouts, the prescribed fence itself suffices. For targets with all pullbacks, a fence of length $L$ extends after padding to length at most $2L$. In the pushout case, every pseudo-Dwyer inclusion has this property.
\end{claim}
\begin{proof}
D.7 handles backward transformations. Suppose first that $F$ is filtered. For a generating cell attachment and a prescribed forward transformation on its boundary, the old cell map together with that transformation is a functor on the finite corner
$P_n\times\{0\}\cup\partial P_n\times[1]$.
Choose a cocone of this finite diagram in $F$. Send every new cell object at the final time to the cone vertex, and use its cone arrows for the new transformation components and the arrows from old boundary objects to new cell objects. The cocone equations are precisely the mixed composition and naturality equations. Thus the one forward transformation extends across the cell without padding. Attach all cells in succession using this same time interval, take the union at limit stages, and restrict along an arrow retract. This proves the filtered case for all Thomason cofibrations.

Now let $F$ have all pushouts, and use the pseudo-Dwyer data $C\subset W\subset D$, $s:W\to C$, and $\varepsilon:is\Rightarrow\mathrm{id}_W$. For a prescribed forward transformation $\alpha:f|_C\Rightarrow h$, define on $W$
\[
 g(w)=f(w)\amalg_{f(is(w))}h(s(w)),
\]
using $f\varepsilon_w$ and $\alpha_{s(w)}$. Choose the pushout on $C$ to be exactly $h(c)$. Universal properties give the functor $g|_W$ and the transformation $f|_W\Rightarrow g|_W$. Outside $W$ retain $f$. There are no arrows from $W$ to its complement, and a mixed arrow from the complement to $W$ is sent through $f(w)\to g(w)$. The same equations prove functoriality and extend the transformation, with its restriction precisely $\alpha$.

Finally suppose that $F$ has all pullbacks. We prove a stronger span-extension assertion: given $f:D\to F$ and a prescribed span $h_0\Rightarrow f|_C$, $h_0\Rightarrow h_1$ on the source of a Thomason cofibration, this span extends over $D$ while its first endpoint stays exactly $f$. For a cell let $A=\partial P_n$. The poset $P_n$ consists of the old copy of $A$, an upper copy of $A$, and a new object $z$ below every upper object. On the upper copy write the given map as $f^+:A\to F$; it receives a transformation from $f|_A$ and a cone from $f(z)$. The prescribed left transformation gives $h_0\to f^+$. Form
\[
 T_a=f(z)\times_{f^+(a)}h_0(a),\qquad a\in A.
\]
These objects and their induced arrows form a finite diagram in $F/f(z)$. That slice has a terminal object and pullbacks, hence all finite limits. Let $L$ be its limit. It has compatible maps to every $h_0(a)$ and to $f(z)$, with exactly the required commutative squares.

The middle cell functor takes its boundary to $h_0$, its upper copy to $h_0$, and its new center to $L$. The last cell functor takes its boundary and upper copy to $h_1$ and its center again to $L$, using the composites $L\to h_0(a)\to h_1(a)$ for the cone. The first transformation goes from the middle functor to $f$, using the specified transformation on the boundary and $L\to f(z)$ at the center. The second uses $h_0\to h_1$ on the boundary and upper copy and $\mathrm{id}_L$ at the center. The limit equations and the two given naturality equations check every arrow. For $n=0$ the same construction uses the terminal object of the slice as the empty limit.

This span construction works cell by cell while keeping both arrows of the span fixed, so it extends over an arbitrary relative cell complex and passes to arrow retracts. For a forward boundary transformation take the first arrow of the boundary span to be the identity; the resulting two-edge fence adds one stationary backward edge before the prescribed forward edge. Apply this separately to each forward edge of the original fence, and D.7 to each backward edge. This proves the bound $2L$.
\end{proof}

\begin{claim}{D.15}{Proved}{Two edges do not suffice for all fibrant targets}
There are an explicitly specified Thomason-fibrant, weakly contractible poset $F$, the generating one-cell cofibration $\{0,4\}\hookrightarrow P_4$, a functor $f:P_4\to F$, and one prescribed forward transformation on its two endpoints for which no extension exists over any padded fence with at most two edges. The full padded extension nevertheless exists by D.10.
\end{claim}
\begin{proof}
Let $K$ be the seven-object poset obtained from $P_4$ by adjoining two objects $y_0,y_1$ and arrows $0\to y_0$, $4\to y_1$. There are no further nonidentity arrows. Define $F_0=K$. For every nonnegative integer $m$, form $F_{m+1}$ by attaching, simultaneously, one generating Thomason horn filler for every horn dimension, horn index, and functor of that horn into $F_m$. Put $F=\operatorname{colim}_{m<\omega}F_m$. This specifies a small poset: each horn map is a cofibration, and a cofibration with poset source has poset target. Finite horn domains make $F$ fibrant. The map $K\to F$ is a trivial Thomason cofibration, so $F$ is weakly contractible because the graph of $K$ is a tree. By \citet[Proposition~1]{D_Cisinski1999}, $K\to F$ is a Dwyer sieve; write $W$ for its cosieve neighborhood and $s:W\to K$ for its right-adjoint retraction.

Let $f$ be the original inclusion of $P_4$, and prescribe the two arrows to $y_0,y_1$ on its endpoints. In $K$, the distance from $0$ to $4$ is four, with a unique shortest path; hence the unique endpoint-preserving functor $P_4\to K$ is $f$. The distance from $y_0$ to $y_1$ is six, so there is no functor $P_4\to K$ with those endpoints. These assertions and every order relation are checked by \texttt{code/d\_two\_step\_obstruction.py}.

Any candidate final functor with endpoints $y_0,y_1$ must have some value outside $W$: otherwise composition with $s$ would give the forbidden functor into $K$. A forward transformation beginning at $f$ has its entire next functor inside $W$, because $W$ is a cosieve containing $K$. A backward transformation ending at $f$ has its entire preceding functor inside $K$, because $K$ is a sieve. When its boundary is stationary at $0,4$, that preceding functor is therefore $f$ itself. The reverse boundary arrows $y_0\to0$ and $y_1\to4$ do not exist, since $K$ is full in $F$.

A fence of at most two edges from the old boundary to the new one has exactly one nonstationary boundary edge. A first forward edge crossing to the new boundary is already impossible because its target lies in $W$. A stationary first forward edge followed by a forward crossing again keeps the final functor in $W$; a backward crossing is prohibited by the absent reverse boundary arrows. A stationary first backward edge leaves the middle functor equal to $f$, so the same argument excludes the second edge. These cases exhaust the orientations and placements of the nonstationary edge. Thus no such extension has at most two edges.
\end{proof}


\subsection{A countable rigid fibrant poset and the unrestricted counterexamples}

For a poset $P$, call $x\in P$ a \emph{downbeat point} when its strict lower set $P_{<x}$ has a greatest element. In particular a minimal element is not a downbeat point. We will use finite principal ideals, rather than any finiteness assumption on the whole poset.

\begin{claim}{D.16}{Proved}{The exact horn-pushout invariant}
Suppose every principal ideal of a poset $P$ is finite and $P$ has no downbeat points. A simultaneous pushout of any family of generating Thomason horn inclusions into $P$ has the same two properties. The inclusion preserves each old principal ideal and preserves the existence or absence of a common upper bound for every pair of old objects.
\end{claim}
\begin{proof}
We give the pushout formula, including the case of noninjective attaching maps. Let $T=\{0,\ldots,n\}$, $M=T\setminus\{i\}$, and
\[
 H=c\operatorname{Sd}^2\Lambda_i^n\ \subset\
 Q=c\operatorname{Sd}^2\Delta^n,
 \qquad n\ge1.
\]
An object $\sigma$ of $Q$ is a nonempty chain of nonempty subsets of $T$, ordered by inclusion of chains. It belongs to $H$ exactly when none of its faces is $M$ or $T$. For a flag $\sigma$, write
$\sigma_H=\sigma\setminus\{M,T\}$.
Whenever this is nonempty, it is an object of $H$. It is the greatest horn subflag of $\sigma$.

For an attaching functor $v:H\to P$, the pushout has the old objects of $P$ and the distinct new flags $Q\setminus H$. The order between new flags is their original inclusion order, there are no arrows from a new flag to an old object, and
\[
 x\le\sigma\quad(x\in P,\ \sigma\notin H)
 \quad\Longleftrightarrow\quad
 \sigma_H\ne\varnothing\ \hbox{ and }\ x\le v(\sigma_H).
 \tag{D.1}
\]
These formulas define a poset and have the pushout universal property. Indeed every old-to-new arrow factors through $v(\sigma_H)$, and the maximal horn-subflag property makes these factorizations compatible with every inclusion of new flags. They are also the explicit Dwyer-pushout formula of \citet[Section~4]{D_Thomason1980}, with the correction of \citet{D_Cisinski1999}. For a simultaneous family, its new objects carry their cell labels and there are no arrows between new objects from different cells; the same formulas apply in each cell.

The old ideals are unchanged. A new flag has finitely many new subflags, and its old predecessors form either the empty set or the single finite ideal $\mathord\downarrow v(\sigma_H)$. Thus all principal ideals remain finite. For two old objects $x,y$, any new common upper bound $\sigma$ gives the old common upper bound $v(\sigma_H)$ by (D.1). The converse implication holds because the old poset is full.

It remains to exclude a new downbeat point. A singleton new flag is minimal, so is not downbeat. Let a new flag $\sigma$ have at least two faces, and choose a new face $N\in\sigma\cap\{M,T\}$. The singleton flag $[N]$ is a new strict predecessor of $\sigma$. A greatest strict predecessor would therefore have to be a new proper subflag $\tau\subsetneq\sigma$ containing $N$, since no new object has an arrow into the old poset. Choose a face $S\in\sigma\setminus\tau$. If $S$ is new, the new singleton $[S]$ is a strict predecessor not below $\tau$. If $S$ is old and $|\sigma|\ge3$, the new two-face subflag $[S,N]$ is a strict predecessor not below $\tau$. If $S$ is old and $|\sigma|=2$, then $\tau=[N]$; the old object $v([S])$ lies below $\sigma$ but cannot lie below the new minimal object $[N]$, by (D.1). Every case contradicts the asserted greatest predecessor. This proof allows every identification made by a noninjective $v$; all subflags used for the first two contradictions are new and remain distinct.
\end{proof}

\begin{claim}{D.17}{Proved}{A countable bifibrant weakly contractible poset with isolated identity}
There is a specified countable poset $F$ which is Thomason-cofibrant and Thomason-fibrant, whose classifying space is contractible, and for which the identity functor is an isolated object of $\mathrm{Fun}(F,F)$. More explicitly, every order-preserving $g:F\to F$ satisfying $g\le\mathrm{id}_F$ or $g\ge\mathrm{id}_F$ is the identity. Consequently its component in $\mathrm{Fun}(F,F)$ consists only of the identity.
\end{claim}
\begin{proof}
Let $F_0=\{a\}$. For each integer $m\ge0$, define $F_{m+1}$ by the simultaneous pushout
\[
\begin{array}{ccc}
\displaystyle\coprod_{n\ge1,\,0\le i\le n,\,v:H_i^n\to F_m} H_i^n
 &\longrightarrow&F_m\\
\big\downarrow&&\big\downarrow\\
\displaystyle\coprod_{n\ge1,\,0\le i\le n,\,v:H_i^n\to F_m} Q^n
 &\longrightarrow&F_{m+1},
\end{array}
\qquad
 H_i^n=c\operatorname{Sd}^2\Lambda_i^n,
 \quad Q^n=c\operatorname{Sd}^2\Delta^n.
\tag{D.2}
\]
Each cell has its full index $(m,n,i,v)$, so its new objects are fresh. In particular \emph{every horn map is attached on every pass, including maps that already admit fillers}. Put $F=\operatorname{colim}_{m<\omega}F_m$ with the full sieve inclusions and the explicit order of (D.1). This specifies the example without choosing any fillers.

There are countably many finite horn shapes, and a finite category has at most countably many functors into a countable poset. Inductively every $F_m$, and hence $F$, is countable. Every stage map is a pushout of generating trivial cofibrations. Therefore $\{a\}\to F$ is a trivial Thomason cofibration, and $F$ is cofibrant and weakly contractible. Every map of a finite horn into $F$ factors through some $F_m$ and receives its designated filler in $F_{m+1}$. Thus $F$ is fibrant. The realization of its nerve is a CW complex, so weak contractibility gives contractibility.

D.16 applies at each pass. Moreover a principal ideal of an object never changes after that object's birth, because all later inclusions are sieves. Hence every principal ideal in $F$ is finite, and $F$ has no downbeat points.

The construction has the following additional branching property:
\[
 x<z\text{ in }F\quad\Longrightarrow\quad
 \text{there is }y>x\text{ with no common upper bound for }y,z.
 \tag{D.3}
\]
Choose a stage $F_m$ containing $x,z$. On the next pass, use the fresh one-horn $\{0\}\hookrightarrow P_4$ attached at $0\mapsto x$. Let $y$ be its new vertex $1$. Within that new cell $y$ is maximal, its old predecessors are exactly $\mathord\downarrow x$, and it has no arrows to the new objects of any other cell in that pass. Thus $y,z$ have no common upper bound in $F_{m+1}$: the only possible upper bound of $y$ there is $y$ itself, while $z\nleq x$ implies $z\nleq y$. D.16 preserves this absence at every later pass. Any proposed upper bound in the union would occur at a finite stage, so none exists in $F$.

Now suppose $g\le\mathrm{id}_F$ and $g$ moves some object. Choose a moved $x$ for which the finite cardinality of $\mathord\downarrow x$ is minimal. Every strict predecessor of $x$ is fixed. Since $g(x)<x$, monotonicity gives
$y=g(y)\le g(x)$ for every $y<x$. Therefore $g(x)$ is the greatest strict predecessor of $x$, contradicting the absence of downbeat points. Thus $g=\mathrm{id}_F$.

Suppose instead that $g\ge\mathrm{id}_F$ and that $g(x)=z>x$. Choose $y$ from (D.3). Then $g(y)\ge y$ and, because $x<y$, monotonicity gives $g(y)\ge g(x)=z$. This is a forbidden common upper bound. Thus again $g=\mathrm{id}_F$. Natural transformations between poset functors are exactly pointwise inequalities. There is consequently no nonidentity natural transformation entering or leaving $\mathrm{id}_F$, and no finite zigzag can leave it.

The program \texttt{code/d\_rigid\_fibrant\_checks.py} verifies the pushout formulas, preservation of old ideals and common-upper-bound relations, and absence of downbeat points for every horn map in dimensions one and two into $\{*\}$ and $P_4$. This includes $2265$ noninjective attaching maps for each two-dimensional horn into $P_4$. It also verifies all constant attaching maps through dimension four and a finite fresh-horn branching pass: $57$ objects after the first pass restricted to dimensions one and two, $285$ after one fresh one-horn at each object, and witnesses for all $140$ strict old pairs. These finite checks certify the local formulas; the preceding induction and finite-horn argument prove the infinite assertions.
\end{proof}

\begin{claim}{D.18}{Refuted}{Unrestricted prescribed-fence homotopy extension}
The prescribed finite-fence homotopy extension property fails for a Thomason trivial cofibration between fibrant posets and a fibrant target. It fails even when the source, target, and codomain of that cofibration have contractible classifying spaces.
\end{claim}
\begin{proof}
Use $F$ from D.17 and $i:\{a\}\hookrightarrow F$, its specified trivial cofibration. Let $b$ be vertex $1$ in the fresh copy of $P_4$ attached at $a$ by the endpoint-$0$ one-horn on the first pass of (D.2). Thus $a<b$ and $a\ne b$. Take $f=\mathrm{id}_F$ and prescribe the single forward transformation $a\to b$ on the one-object source of $i$.

An extension over any finite fence $Z'\to[1]$ preserving the endpoint positions would give a finite zigzag of endofunctors of $F$ from $\mathrm{id}_F$ to a functor $g$ satisfying $g(a)=b$. D.17 forces $g=\mathrm{id}_F$, which instead takes $a$ to $a$. This contradiction excludes every finite length and every allowed orientation or backtracking of the padding fence. The source $\{a\}$ and both occurrences of $F$ are fibrant; the unit $i$ is a trivial cofibration by its explicit construction.
\end{proof}

\begin{claim}{D.19}{Refuted}{The strict-retraction assertion of CATHOT E.6}
There are countable posets $A,B$, with $A$ Thomason-fibrant, a Thomason trivial cofibration $i:A\hookrightarrow B$, and $p:B\to A$ with $pi\simeq\mathrm{id}_A$ by two natural transformations, for which no strict retraction $r:B\to A$ belongs to the categorical homotopy class of $p$. Both classifying spaces are contractible. An explicit strict retraction exists, so the obstruction is exactly its categorical homotopy class.
\end{claim}
\begin{proof}
Keep the specified $F,a,b$ of D.17--D.18. Define $A$ by adjoining to $F$ one new object $c$, with
\[
 c<x\quad\Longleftrightarrow\quad b\le x\text{ in }F,
 \qquad x\in F,
\]
and with no arrows from $F$ to $c$. Let $k:F\hookrightarrow A$, and let $q:A\to F$ be the identity on $F$ with $q(c)=b$. There is a natural transformation $\eta:\mathrm{id}_A\Rightarrow kq$, identity on $F$ and given by $c\to b$ at $c$, and $qk=\mathrm{id}_F$.

The poset $A$ is fibrant. Given a horn map $v:H\to A$, extend $qv$ to its simplex $Q$ using fibrancy of $F$. Keep $v$ on the old horn, and send every new object of $Q$ into the full copy of $F$ using this extension. For a mixed arrow from an old object to a new object, compose the corresponding component of $\eta$ with the extended arrow in $F$. There are no arrows from a new object into the horn because the horn is a sieve. Naturality of $\eta$ and functoriality of the extension check all mixed composites. This is a strict horn filler. Also $qk=\mathrm{id}_F$ and $\eta$ give $BA\simeq BF$, so $BA$ is contractible.

Define $e:A\to A$ by the identity on $F$ and $e(c)=a$. It is a functor because $a<b\le x$ for every mixed relation $c<x$. There is a transformation $e\Rightarrow kq$, identity on $F$ and given by $a\to b$ at $c$. Together with $\eta$ this is
\[
 e\ \Longrightarrow\ kq\ \Longleftarrow\ \mathrm{id}_A.
 \tag{D.4}
\]
Form the specified pushout
\[
 B=A\amalg_{\{c\}=\{a\}}F,
\]
where the second copy of $F$ is attached using the trivial cofibration $\{a\}\hookrightarrow F$ of D.17. Hence $i:A\hookrightarrow B$ is a trivial cofibration. The poset pushout formula makes $B$ a countable poset, and $BB$ is contractible because $i$ is a weak equivalence.

The maps $e:A\to A$ and $k\mathrm{id}_F:F\to A$ agree on the identified object: both take it to the original $a\in F\subset A$. They define $p:B\to A$, and (D.4) proves $pi\simeq\mathrm{id}_A$. For comparison, the maps $\mathrm{id}_A$ on the first summand and the constant functor at $c$ on the second $F$ define an explicit strict retraction $r_0:B\to A$.

Suppose that any strict retraction $r$ satisfies $r\simeq p$ by a finite categorical homotopy. Restrict that homotopy to the second copy of $F$ and compose with $q$. This gives
\[
 g=qr|_F\ \simeq\ qp|_F=\mathrm{id}_F.
\]
D.17 forces $g=\mathrm{id}_F$. But the attaching object $a$ of the second copy is the object $c$ of $A$, which $r$ must fix. Thus
$g(a)=q(r(c))=q(c)=b\ne a$, a contradiction. This proves the required failure with every object, order, and functor specified.
\end{proof}

\begin{claim}{D.20}{Proved}{Failure of comparison even on bifibrant sources and a cylinder obstruction}
For the countable bifibrant poset $F$ of D.17, the maps $\mathrm{id}_F$ and $\operatorname{const}_a$ have the same image in $[BF,BF]$ and distinct classes in $\pi_0\mathrm{Fun}(F,F)$. Moreover the endpoint inclusion
\[
 F\amalg F\longrightarrow F\times P_4
\]
is not a Thomason cofibration, although $F$ is cofibrant and the endpoint inclusion $\{0,4\}\to P_4$ is a generating Thomason cofibration.
\end{claim}
\begin{proof}
Contractibility of $BF$ makes the two realized maps homotopic, while isolatedness of the identity separates their categorical classes. The projection $F\times P_4\to F$ is a weak equivalence, since the finite fence is categorically contractible. If the displayed endpoint inclusion were a cofibration, this product would be a cylinder object for the cofibrant $F$. The two maps into the fibrant $F$ agree in the model homotopy category, so independence of cylinder object would produce a homotopy through this cylinder \citep[Proposition~A.2.2.1]{CpathsLurieHTT}. That would be a four-edge categorical homotopy between the two maps, contradicting D.17. Thus the endpoint inclusion fails the cofibration condition.
\end{proof}

\paragraph{Conclusion for Part D.}
D.18 refutes the unrestricted prescribed-fence extension property, and D.19 refutes the additional-fibrancy strict-retraction assertion of CATHOT E.6, including the case of a trivial cofibration and contractible classifying spaces. The requested smallest single-cell instance is affirmative by D.10. D.12--D.13 prove all finite relative cell and finite-presentation cases; the countable examples above show why no unrestricted passage from those cases to an infinite cell complex is valid. The weaker Hot-class retraction in D.1 remains true and is realized explicitly by $r_0$ in D.19, whose realized map has the same class as $p$.
